% ===================== 2. CẤU TRÚC DỮ LIỆU & TRẠNG THÁI =====================
\section{Cấu trúc dữ liệu và quản lý trạng thái (State Representation)}
\label{sec:data}

Hiệu năng của cả ba tầng phụ thuộc vào việc \emph{kiểm tra hợp lệ} và \emph{cập
nhật} một phép gán phải thật nhanh, vì các thuật toán xấp xỉ thực hiện hàng triệu
thao tác này. Mục này trình bày biểu diễn dữ liệu được thiết kế cho mục tiêu đó.

\subsection{Biểu diễn đối tượng (ClassDTO, RoomDTO)}
Hai đối tượng dữ liệu thuần (DTO) tách \emph{thuộc tính đầu vào} khỏi \emph{trạng
thái lời giải}:
\begin{itemize}
  \item \texttt{ClassDTO}: bất biến $(t_i,g_i,s_i)$ cộng phần trạng thái thay đổi
  theo lời giải --- \texttt{assigned\_slot} (tiết bắt đầu, 1-based), \texttt{assigned\_room},
  \texttt{is\_assigned}. Nhờ tách bạch này, mọi thuật toán dùng \emph{chung} một
  biểu diễn lớp; chỉ phần trạng thái bị ghi đè khi gán/gỡ.
  \item \texttt{RoomDTO}: định danh phòng và sức chứa $c_v$.
\end{itemize}
Các mảng \texttt{classes[\,]}, \texttt{rooms[\,]} dùng chỉ số \emph{1-based} (phần
tử $0$ để trống) để khớp định danh lớp/phòng trong đề và định dạng xuất.

\subsection{Ma trận trạng thái thời gian (room\_busy, teacher\_busy)}
Trái tim của lớp \texttt{ScheduleState} là hai ma trận bận theo thời gian:
\begin{center}
\begin{tabular}{ll}
\toprule
\texttt{room\_busy[$r$][$\tau$]} & phòng $r$ có bị chiếm tại tiết $\tau$ không \\
\texttt{teacher\_busy[$g$][$\tau$]} & giáo viên $g$ có bận tại tiết $\tau$ không \\
\bottomrule
\end{tabular}
\end{center}
Đây là hiện thực trực tiếp của hai họ ràng buộc xung đột~\eqref{eq:c3}--\eqref{eq:c4}:
một phép gán hợp lệ tương đương với việc cả khối $t_i$ tiết của lớp đều \emph{rảnh}
trên cả hai ma trận. \texttt{teacher\_busy} dùng \texttt{dict} (chỉ sinh hàng cho
giáo viên thực sự xuất hiện) để tiết kiệm bộ nhớ khi mã giáo viên thưa.

\begin{figure}[H]
\centering
\begin{tikzpicture}[
  cell/.style={draw, minimum size=0.55cm},
  busy/.style={cell, fill=accent!40},
  free/.style={cell, fill=white},
  target/.style={cell, fill=cfast!30, draw=cfast, very thick},
  rowlbl/.style={font=\scriptsize, anchor=east},
  collbl/.style={font=\scriptsize, anchor=south},
  tcol/.style={font=\scriptsize\bfseries, anchor=south, text=cfast!75!black},
  gtitle/.style={font=\small\bfseries, anchor=south}
]
  % --- annotation chung (đặt trên cùng, không đè tiêu đề) ---
  \node[font=\scriptsize\itshape, text=cfast!75!black, align=center] at (3.6,1.75)
    {Khối ứng viên $t_c{=}3$ tiết: cần cùng \emph{trống} tại $\tau_2,\tau_3,\tau_4$ trên \textbf{cả hai} ma trận};

  % --- room_busy ---
  \node[gtitle] at (1.375,0.95) {\texttt{room\_busy}};
  \foreach \x in {0,1,5} {\node[collbl] at (\x*0.55,0.16){$\tau_{\x}$};}
  \foreach \x in {2,3,4} {\node[tcol] at (\x*0.55,0.16){$\tau_{\x}$};}
  \node[rowlbl] at (-0.35,-0.275){$r_1$};
  \node[rowlbl] at (-0.35,-0.825){$r_2$};
  \node[rowlbl] at (-0.35,-1.375){$r_3$};
  \foreach \x/\s in {0/busy,1/busy,2/free,3/free,4/free,5/busy} {\node[\s] at (\x*0.55,-0.275){};}
  \foreach \x/\s in {0/free,1/free,2/target,3/target,4/target,5/free} {\node[\s] at (\x*0.55,-0.825){};}
  \foreach \x/\s in {0/busy,1/free,2/busy,3/busy,4/free,5/free} {\node[\s] at (\x*0.55,-1.375){};}

  % --- teacher_busy ---
  \begin{scope}[xshift=5.0cm]
    \node[gtitle] at (1.375,0.95) {\texttt{teacher\_busy}};
    \foreach \x in {0,1,5} {\node[collbl] at (\x*0.55,0.16){$\tau_{\x}$};}
    \foreach \x in {2,3,4} {\node[tcol] at (\x*0.55,0.16){$\tau_{\x}$};}
    \node[rowlbl] at (-0.35,-0.275){$g_1$};
    \node[rowlbl] at (-0.35,-0.825){$g_2$};
    \foreach \x/\s in {0/free,1/busy,2/free,3/free,4/busy,5/busy} {\node[\s] at (\x*0.55,-0.275){};}
    \foreach \x/\s in {0/free,1/free,2/target,3/target,4/target,5/free} {\node[\s] at (\x*0.55,-0.825){};}
  \end{scope}

  % --- chú giải ---
  \node[busy] at (0,-2.05){}; \node[anchor=west,font=\scriptsize] at (0.35,-2.05){bận};
  \node[free] at (1.7,-2.05){}; \node[anchor=west,font=\scriptsize] at (2.05,-2.05){trống};
  \node[target] at (3.5,-2.05){}; \node[anchor=west,font=\scriptsize] at (3.85,-2.05){khối ứng viên};
\end{tikzpicture}
\caption{Kiểm tra tính hợp lệ: để xếp một lớp $t_c{=}3$ tiết vào phòng $r_2$ do
giáo viên $g_2$ dạy, hệ thống chỉ cần tìm một khối $3$ ô \emph{cùng trống} ở cùng
các tiết ($\tau_2,\tau_3,\tau_4$) trên \texttt{room\_busy}[$r_2$] \textbf{và}
\texttt{teacher\_busy}[$g_2$].}
\label{fig:state-matrices}
\end{figure}

\paragraph{Chèn / gỡ trong $\mathcal{O}(t_i)$.} Thao tác \texttt{insert} đánh dấu
$t_i$ ô bận trên hai ma trận và ghi trạng thái vào \texttt{ClassDTO}; \texttt{remove}
làm ngược lại. Cả hai chỉ tốn $\mathcal{O}(t_i)\le\mathcal{O}(4)$ --- gần như hằng
số. Ngoài ra \texttt{ScheduleState} cung cấp \texttt{snapshot}/\texttt{restore}
nhẹ (chỉ lưu ánh xạ \texttt{class\_id}$\to$(phòng, tiết)) để metaheuristic có thể
thử một nước đi rồi hoàn tác tức thì --- nền tảng cho Local Search và ALNS.

\subsection{Đánh giá hợp lệ cục bộ (Fast Feasibility Check)}
Hàm \texttt{can\_place($c,r,s_0$)} trả lời ``lớp $c$ có đặt được vào phòng $r$ từ
tiết $s_0$ không'' theo đúng thứ tự cắt tỉa từ rẻ tới đắt:
\begin{enumerate}[leftmargin=*]
  \item Kiểm tra sức chứa $c_r \ge s_c$ (một phép so sánh) --- loại nhanh phần lớn
  phòng không hợp lệ.
  \item Quét $t_c$ ô trên \texttt{room\_busy} và \texttt{teacher\_busy}; gặp ô bận
  đầu tiên là trả \texttt{False} ngay (short-circuit).
\end{enumerate}
Cùng nhóm tiện ích: \texttt{find\_position} (vị trí trống đầu tiên cho một lớp) và
\texttt{feasible\_room\_count} (đếm số phòng khả thi --- dùng cho Regret-2 ở
Mục~\ref{sec:alns}). Tập $\mathcal{S}(t_i)$ được \emph{cache} theo thời lượng để
không phải sinh lại. Nhờ thiết kế này, một lần kiểm tra hợp lệ chỉ tốn
$\mathcal{O}(t_i)$, cho phép các vòng lặp tìm kiếm chạy hàng chục nghìn lần/giây.
